\par\Needspace{6\baselineskip}\medskip
\hypertarget{native-117-article087}{}\pdfbookmark[2]{87.}{article-087}
\begin{center}87.\end{center}
Die folgenden Beispiele veranschaulichen die Anwendung der vorhergehenden Rechenvorschriften und lassen zugleich den Grad der Genauigkeit abschätzen.

I. Es sei $\log r=0{,}330\,7640$, $\log r'=0{,}322\,2239$, $\Delta=7^0\,34'\,53''{,}73=27293''{,}73$ und $t=21{,}93391$ Tage. Daraus erhält man $\omega=-33'\,47''{,}90$. Die weitere Rechnung verläuft wie folgt.
\SourceLeaf{118}
\begin{center}
\begin{tabular}{@{}ll@{}}
\makebox[17em][l]{$\log\Delta$\hfill .\,.\,.\,.\,.\,.\,.\,.\,.\,.\,.} & $4{,}436\,0629$\\[.3em]
\makebox[17em][l]{$\log rr'$\hfill .\,.\,.\,.\,.\,.\,.\,.\,.\,.} & $0{,}652\,9879$\\[.3em]
\makebox[17em][l]{$\mathrm{C.}\log3k$\hfill .\,.\,.\,.\,.\,.\,.\,.} & $5{,}972\,8722$\\[.3em]
\makebox[17em][l]{$\mathrm{C.}\log t$\hfill .\,.\,.\,.\,.\,.\,.\,.\,.} & $8{,}658\,8840$\\[.3em]
\makebox[17em][l]{$\mathrm{C.}\log\cos2\omega$\hfill .\,.\,.\,.} & $0{,}000\,0840$\\[.3em]
\hline
\makebox[17em][l]{$\log a$\hfill .\,.\,.\,.\,.\,.\,.\,.\,.\,.\,.} & $9{,}720\,8910$\\[.3em]
\end{tabular}
\end{center}
\begin{center}
\begin{tabular}{@{}ll@{}}
\makebox[17em][l]{$\tfrac12\log rr'\cos2\omega$\hfill .\,.\,.} & $0{,}326\,4519$\\[.3em]
\makebox[17em][l]{$2\log\sin\tfrac14\Delta$\hfill .\,.\,.\,.\,.} & $7{,}038\,9972$\\[.3em]
\makebox[17em][l]{$\log\tfrac{2}{27}$\hfill .\,.\,.\,.\,.\,.\,.\,.\,.\,.} & $8{,}869\,6662$\\[.3em]
\makebox[17em][l]{$\mathrm{C.}\log aa$\hfill .\,.\,.\,.\,.\,.\,.} & $0{,}558\,2180$\\[.3em]
\makebox[17em][l]{$\mathrm{C.}\log\cos\omega$\hfill .\,.\,.\,.\,.\,.} & $0{,}000\,0210$\\[.3em]
\hline
\makebox[17em][l]{$\log\beta$\hfill .\,.\,.\,.\,.\,.\,.\,.\,.\,.\,.} & $6{,}793\,3543$\\[.3em]
$\beta=$ & $0{,}000\,6213\,757$\\[.3em]
\end{tabular}
\end{center}
\begin{center}
\begin{tabular}{@{}ll@{}}
\makebox[17em][l]{$\log2$\hfill .\,.\,.\,.\,.\,.\,.\,.\,.\,.\,.} & $0{,}301\,0300$\\[.3em]
\makebox[17em][l]{$2\log\cos\tfrac12\Delta$\hfill .\,.\,.\,.\,.} & $9{,}998\,0976$\\[.3em]
\makebox[17em][l]{$2\log\cos2\omega$\hfill .\,,\,.\,.\,.} & $9{,}999\,8320$\\[.3em]
\makebox[17em][l]{$\mathrm{C.}\log(1-3\beta)$\hfill .\,.\,.\,.} & $0{,}000\,8103$\\[.3em]
\makebox[17em][l]{$2\mathrm{C.}\log\cos\omega$\hfill .\,.\,.\,.} & $0{,}000\,0420$\\[.3em]
\hline
\makebox[17em][l]{$\log\gamma$\hfill .\,.\,.\,.\,.\,.\,.\,.\,.\,.\,.} & $0{,}299\,8119$\\[.3em]
$\gamma=$ & $1{,}994\,3982$\\[.3em]
$21\beta=$ & $0{,}013\,0489$\\[.3em]
\end{tabular}
\end{center}
\begin{center}
\begin{tabular}{@{}ll@{}}
$1+\gamma+21\beta=$ & $3{,}007\,4471$\\[.3em]
\hline
\makebox[17em][l]{$\log$\hfill .\,.\,.\,.\,.\,.\,.\,.\,.\,.\,.\,.} & $0{,}478\,1980$\\[.3em]
\makebox[17em][l]{$\log a$\hfill .\,.\,.\,.\,.\,.\,.\,.\,.\,.\,.} & $9{,}720\,8910$\\[.3em]
\makebox[17em][l]{$\mathrm{C.}\log(1+5\beta)$\hfill .\,.\,.\,.} & $9{,}998\,6528$\\[.3em]
\hline
\makebox[17em][l]{$\log\BareRoot{p}$\hfill .\,.\,.\,.\,.\,.\,.\,.\,.} & $0{,}197\,7418$\\[.3em]
\makebox[17em][l]{$\log p$\hfill .\,.\,.\,.\,.\,.\,.\,.\,.\,.\,.} & $0{,}395\,4836$\\[.3em]
\end{tabular}
\end{center}
\noindent Dieser Wert von $\log p$ weicht vom wahren Wert um kaum eine Einheit an der siebten Dezimalstelle ab. Formel VI ergibt in diesem Beispiel $\log p=0{,}395\,4822$, Formel VII ergibt $0{,}395\,4780$ und Formel VIII schließlich $0{,}395\,4754$.

II. Es sei $\log r=0{,}428\,2792$, $\log r'=0{,}406\,2033$, $\Delta=62^0\,55'\,16''{,}64$ und $t=259{,}88477$ Tage. Daraus erhält man $\omega=-1^0\,27'\,20''{,}14$, $\log a=9{,}748\,2348$, $\beta=0{,}045\,3511\,4$, $\gamma=1{,}681\,121$, $\log\BareRoot{p}=0{,}219\,8013$ und $\log p=0{,}439\,6026$. Der letzte Wert ist um $211$ Einheiten an der siebten Dezimalstelle zu klein. Der wahre Wert in diesem Beispiel ist nämlich $0{,}439\,6237$. Formel VI ergibt $0{,}436\,8730$, Formel VII ergibt $0{,}415\,9824$ und Formel VIII schließlich $0{,}405\,1103$. Die beiden letzten Werte weichen hier so stark vom wahren Wert ab, dass sie nicht einmal als Näherungen dienen können.
